DGFiP, IRCOM revenus 2024 ; calculs des auteurs.\par Note : taux moyen = IR et CSG payés sur les revenus imposés au barème (activité, pensions, indépendants, foncier), rapportés à leur montant brut (assiette de type CSG), par foyer fiscal. La bande grise encadre le taux actuel selon le partage de l'IR entre ces revenus et ceux qui restent au PFU (12,8 \% de ces derniers, ou prorata). Barèmes : A, 13,3 \% dès le premier euro ; C, huit tranches de 5 \% à 40 \% ; E, 0 \% jusqu'à 14 000 \euro{}, 17,2 \% jusqu'à 120 000 \euro{}, 40 \% jusqu'à 250 000 \euro{}, 49 \% au-delà.
